\scanpage{76}
\mspage{67}
% archive sheet number: 76

travaille avec la relation $X\mathrel{\|}Y\Longleftrightarrow\widetilde{X}\cap\widetilde{Y}=\varnothing$, les supports sont \uline{certaines} parties fermées de $M,\leq$ bien particulières. \unsure{Jamais par} elles on ne pourra exprimer des supports internes !

Je désigne par $\Sigma_M$ (notation \unsure{standard|standart}) \unsure{l'ens} des supports pour $M$. \deleted{L'application} Cet ensemble, \deleted{\unsure{est|et}} ordonné par $\subset$, a un plus petit élément $\varnothing$, et un plus grand élément $M$. Il a des $\operatorname{Inf}$ quelconques -- ce sont les $\bigcap_i$ simplement. Il y a une \deleted{fonction} application
\[
\Sigma_M\overset{\sim}{\longrightarrow}\Sigma_M\qquad S\longrightarrow\deleted{\unsure{C}}\unsure{\complement|C}S,
\]
qui est une antiinvolution de l'ens ordonné. Il \deleted{\unsure{tr}} transforme des $\operatorname{Inf}$ en $\operatorname{Sup}$ et inversement. Le $\operatorname{Sup}$ s'écrit aussi
\[
\operatorname{Sup}S_i=\operatorname{supp}^{\circ}(\textstyle\bigcup S_i)
\]
On prendra bien garde de ne pas \unsure{le|les} confondre avec $\bigcup S_i$.

On peut aussi construire un $\Sigma_L$, avec $\mid\circ\mid_L=\mathrel{\|}_L$ sur $L$. On \deleted{\unsure{trouvera}}\inserted{voudrait} une application canonique
